\section{Modelo conceptual}
La Figura~\ref{fig:conceptual} presenta las entidades, relaciones y cardinalidades del dominio.
\begin{figure}[ht]
\centering
\includegraphics[width=\textwidth]{modelo_conceptual.png}
\caption{Modelo conceptual del dominio.}
\label{fig:conceptual}
\end{figure}

El modelo separa cuatro núcleos relacionados. El núcleo operativo reúne
productos, clientes, condiciones comerciales, pedidos y entregas. El núcleo
documental organiza documentos, versiones y fragmentos. El acceso vincula cada
actor con un único perfil y cada perfil con las clases documentales permitidas.
La interacción conserva consultas, respuestas, evidencias y eventos de
auditoría.

Las cardinalidades expresan reglas del dominio. Un pedido pertenece a un
cliente y contiene una o más líneas; una entrega pertenece a un pedido, que
puede tener entregas parciales; una incidencia pertenece a una entrega. Un
documento puede tener varias versiones y cada versión se divide en fragmentos
ordenados. Cada fragmento puede conservar un embedding por versión del modelo.
Una respuesta exitosa se relaciona con al menos una evidencia documental o
estructurada, mientras que los resultados negativos no admiten evidencia.

El historial forma parte del modelo conceptual. Las condiciones comerciales y
las versiones documentales tienen vigencia temporal; una versión sustituida o
revocada deja de participar en la recuperación, pero no se elimina porque
puede haber sostenido una respuesta anterior.
